\section{Informationsfluss über die Attention}

Dieses Kapitel untersucht, welche \glspl{attentionhead} ihre Aufmerksamkeit auf sentimenttragende Wörter richten, und damit als Kandidaten für die kausale Verarbeitung in Hypothese H1 in Frage kommen. Jeder Kopf berechnet für jede Position eine gewichtete Summe über alle vorherigen Token; ein Gewicht nahe 1 bedeutet, dass eine Position fast ihre gesamte Aufmerksamkeit auf ein einzelnes anderes Token legt. Man kann sich das als eine Frage-und-Antwort-Beziehung vorstellen: Jedes Token stellt über seine Query eine Art Suchanfrage, und jedes andere Token bietet über seinen Key eine passende oder unpassende Antwort an. Je besser Anfrage und Antwort zusammenpassen, desto höher das Attention-Gewicht und desto mehr Information fließt von diesem Token. Bei Pythia-410M gibt es 24 Schichten mit je 16 solchen Köpfen, also 384 Stellen, an denen Information auf diese Weise fließen kann.

An einem Beispielsatz wird untersucht, welche Köpfe besonders stark auf das Wort „delicious“ achten. Die stärkste gemessene Aufmerksamkeit liegt in Schicht 6, Kopf 3, mit einem Gewicht von 0,835 (Abbildung~\ref{fig:subject}); weitere auffällige Köpfe liegen in den Schichten 0 bis 12. Die stärksten Muster konzentrieren sich also auf frühe bis mittlere Schichten, genau dort, wo der Logit Lens noch kaum eine Sentiment-Trennung zeigt.

\begin{figure}[htbp]
  \centering
  \includegraphics[width=0.7\textwidth]{attention_subject_heatmap.png}
  \caption{Aufmerksamkeit der Köpfe auf das sentimenttragende Token „delicious“ an einem Beispielsatz.}
  \label{fig:subject}
\end{figure}

Ein einzelner Beispielsatz kann jedoch täuschen. Aggregiert über alle Prompts des CAD-Datensatzes relativiert sich dieser Befund deutlich: Der höchste gemittelte Aufmerksamkeitswert auf sentimenttragende Wörter liegt bei nur 0,019 (Abbildung~\ref{fig:cadattn}), fast zwei Größenordnungen niedriger als am Einzelbeispiel, und Schicht 6/Kopf 3 fällt im Aggregat nicht besonders auf. Systematisch über den Datensatz hinweg ist die Aufmerksamkeit auf einzelne Sentiment-Wörter also schwach und diffus über viele Köpfe verteilt, statt auf wenige dominante Köpfe konzentriert zu sein.

\begin{figure}[htbp]
  \centering
  \includegraphics[width=0.8\textwidth]{attention_cad_sentiment_heatmap.png}
  \caption{Durchschnittliche Aufmerksamkeit je Schicht und Kopf auf sentimenttragende Wörter über den CAD-Datensatz (n=3413 Prompts).}
  \label{fig:cadattn}
\end{figure}

Ergänzend wurde nach Induction Heads gesucht, also Köpfen, die Muster der Form „A B ... A“ zu „B“ vervollständigen und als bekannter Mechanismus zur Mustererkennung gelten. Die höchsten Induction-Scores (zwischen 0,38 und 0,44) liegen in den frühen Schichten 1, 3, 4 und 5, also in einem ähnlichen Bereich wie die stärkste sentimentbezogene Aufmerksamkeit, aber nicht bei denselben Köpfen.

Damit lässt sich nur mit Vorbehalt festhalten: Es gibt Köpfe, die auf sentimenttragende Wörter fokussieren, aber dieser Effekt ist im Aggregat schwach. Vor allem zeigt hohe Aufmerksamkeit nur, dass ein Kopf Information von einem Token \emph{bezieht}, nicht, dass diese Information für die Vorhersage tatsächlich \emph{gebraucht} wird. Attention beschreibt damit nur den Informationsfluss, nicht die Kausalität. Ob ein Kopf oder eine Schicht wirklich notwendig ist, kann nur ein aktiver Eingriff zeigen; das leistet das Activation Patching im folgenden Kapitel.
